\subsection*{一步一步建立直观}

\paragraph{1. 把消失的量看成一阶位移}

当两个函数同时趋零时，它们的比值可能仍有极限。若在目标点附近分别近似为 $f'(a)h$ 和 $g'(a)h$，共同的小步长可以消去，留下局部变化率之比。选择正弦函数，观察原点附近正弦曲线与斜率为一的直线如何贴合。



\paragraph{2. 先弄清为什么允许比较导数}

窗口切换到平方函数，展示差商的控制机制，但它本身不是一般洛必达法则的证明。一般法则还需要区间上的可导性、分母导数不为零及导数比值的极限等条件。其理由来自中值定理对邻域中变化率的联系，不能只靠画两条切线代替。



\paragraph{3. 若一阶项消失，检查更高阶结构}

改用立方函数并把观察点移到零，一阶导数也为零。此时分子与分母可能以不同阶数消失，必须继续分析主导项，或在每次确认条件后再次使用法则。不要把任意分式直接替换成导数之比；只有符合条件的极限问题才有这种结论。
